\documentclass[11pt,reqno]{amsart}

% --- Encoding / fonts ---
\usepackage[T1]{fontenc}
\usepackage[utf8]{inputenc}

% --- Core AMS packages ---
\usepackage{amsmath,amssymb,amsfonts,amsthm,mathtools}

% --- Graphics and lists ---
\usepackage{graphicx}
\usepackage{enumitem}

% --- Page layout ---
\usepackage[a4paper,margin=3cm]{geometry}

% --- Theorem environments ---
\theoremstyle{plain}
\newtheorem{theorem}{Theorem}[section]
\newtheorem{lemma}[theorem]{Lemma}
\newtheorem{proposition}[theorem]{Proposition}
\newtheorem{corollary}[theorem]{Corollary}

\theoremstyle{definition}
\newtheorem{definition}[theorem]{Definition}
\newtheorem{example}[theorem]{Example}
\newtheorem{question}[theorem]{Question}
\newtheorem{conjecture}[theorem]{Conjecture}

\theoremstyle{remark}
\newtheorem{remark}[theorem]{Remark}
\newtheorem{claim}[theorem]{Claim}
\newtheorem{assertion}[theorem]{Assertion}
\newtheorem{notations}[theorem]{Notations}

\newtheorem{assumption}{Assumption}
\newtheorem{hypothesis}{Hypothesis}

\newtheorem*{theoremA}{Theorem A}
\newtheorem*{theoremB}{Theorem B}
\newtheorem*{theoremC}{Theorem C}

% --- Operators ---
\DeclareMathOperator{\IFS}{IFS}
\DeclareMathOperator{\diam}{diam}
\DeclareMathOperator{\vol}{vol}

% --- Hyperref: load late ---
\usepackage[pdfencoding=auto,psdextra]{hyperref}

% --- TOC depth ---
\setcounter{tocdepth}{3}

% --- Title & Authors ---
\title[T]{TransTitle}

\author[A1]{B1}
\address{Department of Pure Mathematics, Ferdowsi University of Mashhad, Mashhad, Iran}
\email{htina@um.ac.ir}

\author[A2]{B2}
\address{Department of Pure Mathematics, Ferdowsi University of Mashhad, Mashhad, Iran}
\email{zaj@ferdowsi.um.ac.ir}

\thanks{Corresponding author: F. H. Ghane (\texttt{htina@um.ac.ir})}

\keywords{Monotone}
\subjclass[2020]{Primary 37D25; Secondary 37A35, 37C40, 37A60.}

\begin{document}

\maketitle


\begin{abstract}
abbbaaaaaaaa
\end{abstract}

\section{Introduction}
\label{sec:intro}


\section{Model}
\label{sec:model}

\section{Introduction and model formulation}
\label{sec:introduction}

\subsection{Physical origin and governing equations}

The transition from laminar to turbulent pipe flow has been investigated by Barkley and collaborators using a reduced model coupling turbulence intensity to local centerline velocity. The original program was initiated in \cite{barkley2011simplifying} and developed into the model used to describe the rise of fully turbulent flow in \cite{barkley2015rise, barkley2016theoretical,frishman2022dynamical}.

The spatially homogeneous reaction dynamics associated with the later model are
\begin{equation}
\begin{aligned}
\dot q &=q\bigl[r+u-2-(r+0.1)(q-1)^2\bigr],\\
\dot u &=\varepsilon\bigl[2-u+2q(1-u)\bigr],
\end{aligned}
\label{eq:barkley}
\end{equation}
where $q\ge0$ represents turbulence intensity, $u$ is the centerline-velocity variable, $r>0$ is a Reynolds-number-related model parameter, and $\varepsilon>0$ is a timescale parameter. Equations~\eqref{eq:barkley} arise by suppressing the spatial-derivative terms in the published reaction--diffusion--advection model; they are not a new artificial modification \cite{barkley2015rise}.

\section{Model formulation}
\label{sec:model}

Barkley's reduced description of transitional pipe flow couples
turbulence intensity to the local centerline velocity
\cite{barkley2011simplifying,barkley2015rise,barkley2016theoretical}.
Suppressing the spatial-derivative terms gives the homogeneous
reaction system
\begin{equation}
\begin{aligned}
\dot q
&=q\left[r+u-2-(r+\delta)(q-1)^2\right],\\
\dot u
&=\varepsilon\left[2-u+\kappa q(1-u)\right],
\end{aligned}
\qquad
r,\delta,\varepsilon,\kappa>0.
\label{eq:barkley-general}
\end{equation}
Here $q\geq0$ measures turbulence intensity, $u$ denotes the
normalized centerline velocity, and $r$ is a Reynolds-number-related
control parameter. The velocity normalization assigns the value
$u=2$ to the laminar state, while the turbulence-dependent term
drives $u$ toward the reference level $u=1$.
The parameter $\varepsilon$ controls the relative timescale of
velocity relaxation, and $\kappa$ measures the strength of its
coupling to turbulence.

The parameter $\delta$ controls the linear decay of turbulence
near the laminar equilibrium $(q,u)=(0,2)$: the corresponding
eigenvalues are $-\delta$ and $-\varepsilon$. Thus $\delta>0$
preserves the local stability of the laminar state, while
$\kappa>0$ represents the reduction of centerline velocity
associated with increasing turbulence.

In this investigation, we adopt the parameter choices
$\delta=0.1$ and $\kappa=2$, obtaining
\begin{equation}
\begin{aligned}
\dot q
&=q\bigl[r+u-2-(r+0.1)(q-1)^2\bigr],\\
\dot u
&=\varepsilon\bigl[2-u+2q(1-u)\bigr],
\end{aligned}
\qquad r,\varepsilon>0.
\label{eq:barkley}
\end{equation}
These choices fix the laminar decay rate and the
turbulence--velocity coupling strength, leaving $r$ and
$\varepsilon$ as the bifurcation parameters.
System~\eqref{eq:barkley} is a cubic polynomial differential
system. Our aim is to investigate the existence, stability,
and number of its limit cycles in the physical half-plane
$q\geq0$. Such cycles describe spatially homogeneous temporal
oscillations of the reduced model; they should be distinguished
from traveling or spatially localized structures of the
full pipe-flow equations.

The nonzero equilibria satisfy
\[
2(r+0.1)q^3-3(r+0.1)q^2+2(1-r)q+0.1=0.
\]
Dividing by $2(r+0.1)$ and substituting $q=z+\tfrac12$
reduces this algebraic equation to
\[
z^3+p(r)z+b(r)=0,
\]
where
\[
p(r)=\frac{37-70r}{40(r+0.1)},
\qquad
b(r)=\frac{21-30r}{40(r+0.1)}.
\]
This substitution is used only to solve the equilibrium
equation; it does not change the variables of the
differential system.

For $r>r_c\approx0.645189835$, define
\begin{equation}
s(r)=\sqrt{\frac{70r-37}{120(r+0.1)}},
\qquad
\theta(r)=\frac13\arccos\left(
\frac{30r-21}{80(r+0.1)s(r)^3}
\right).
\label{eq:cubic-auxiliary}
\end{equation}
Here $\arccos$ takes its principal value in $[0,\pi]$.
The two positive roots are explicitly
\begin{equation}
\boxed{
\begin{aligned}
q_-(r)
&=\frac12+2s(r)
\cos\left(\theta(r)-\frac{2\pi}{3}\right),\\
q_+(r)
&=\frac12+2s(r)\cos\theta(r).
\end{aligned}}
\label{eq:explicit-positive-roots}
\end{equation}
They satisfy
\[
0<q_-(r)<q_+(r)<2.
\]
Consequently, the two positive equilibria are
\begin{equation}
\boxed{
E_\pm(r)=
\left(
q_\pm(r),
\frac{2[1+q_\pm(r)]}{1+2q_\pm(r)}
\right).
}
\label{eq:explicit-positive-equilibria}
\end{equation}

The remaining root is
\[
q_{\mathrm n}(r)=
\frac12+2s(r)
\cos\left(\theta(r)+\frac{2\pi}{3}\right),
\]
and satisfies $-\tfrac12<q_{\mathrm n}(r)<0$.
Its equilibrium is outside the physical half-plane.

At $r=r_c$, the same formulas hold by continuity:
$\theta(r_c)=\pi/3$ and
\[
q_-(r_c)=q_+(r_c)=q_c.
\]
For $0<r<r_c$, there are no positive roots.

For example, at $r=0.65$ these exact expressions give
\[
\begin{aligned}
E_-&\approx(0.708427852,\;1.413760738),\\
E_+&\approx(0.896537216,\;1.358028411).
\end{aligned}
\]
Together with $E_0=(0,2)$, these are all physical
equilibria for this parameter value.












ooooooooooooooooooooooooo

\subsection{Equilibria and conditions for limit-cycle existence}

Consider
\begin{equation}
\begin{aligned}
\dot q&=q\left[r+u-2-(r+0.1)(q-1)^2\right],\\
\dot u&=\varepsilon\left[2-u+2q(1-u)\right],
\end{aligned}
\qquad r,\varepsilon>0.
\label{eq:model}
\end{equation}

\paragraph{Equilibria.}
Setting $\dot q=\dot u=0$ gives the laminar equilibrium
\[
E_0=(0,2).
\]
Every other equilibrium has the form
\begin{equation}
E_*=\left(q_*,\frac{2(1+q_*)}{1+2q_*}\right),
\label{eq:nonzero-equilibrium}
\end{equation}
where $q_*$ is a real root of
\begin{equation}
2(r+0.1)q^3-3(r+0.1)q^2
+2(1-r)q+0.1=0.
\label{eq:equilibrium-cubic}
\end{equation}
Thus \eqref{eq:equilibrium-cubic}, together with
\eqref{eq:nonzero-equilibrium}, specifies all nonzero
equilibria exactly.

There is always exactly one negative root
$q_{\mathrm n}\in(-1/2,0)$. Its equilibrium lies outside
the physical region $q\geq0$.

To count the positive roots, let $q_c$ be the unique
positive root of
\[
44q_c^3-30q_c^2-3q_c-1=0,
\]
and define
\[
r_c=
\frac{1+22q_c}{10q_c(2-q_c)(1+2q_c)}-\frac1{10}.
\]
Numerically,
\[
q_c\approx0.802140182,
\qquad r_c\approx0.645189835.
\]
The physical equilibria are therefore:
\[
\begin{array}{c|c}
\text{Parameter range} & \text{Equilibria in }q\geq0\\ \hline
0<r<r_c & E_0\\
r=r_c & E_0,\ E_c\text{ (a double positive equilibrium)}\\
r>r_c & E_0,\ E_-,\ E_+
\end{array}
\]
where
\[
E_\pm=
\left(q_\pm,\frac{2(1+q_\pm)}{1+2q_\pm}\right),
\qquad 0<q_-<q_+<2.
\]
The equilibrium coordinates do not depend on $\varepsilon$.

\paragraph{Linear stability.}
At $E_0$, the eigenvalues are
\[
\lambda_1=-0.1,\qquad \lambda_2=-\varepsilon,
\]
so the laminar equilibrium is always locally
asymptotically stable.

At a nonzero equilibrium, the Jacobian is
\begin{equation}
J_*=
\begin{pmatrix}
2(r+0.1)q_*(1-q_*) & q_*\\[1mm]
-\dfrac{2\varepsilon}{1+2q_*}
&-\varepsilon(1+2q_*)
\end{pmatrix}.
\label{eq:jacobian}
\end{equation}
Its trace and determinant are
\begin{align}
T_*&=
2(r+0.1)q_*(1-q_*)-\varepsilon(1+2q_*),
\label{eq:trace}\\
\Delta_*&=
2\varepsilon q_*
\left[
(r+0.1)(q_*-1)(1+2q_*)
+\frac1{1+2q_*}
\right].
\label{eq:determinant}
\end{align}
The negative equilibrium and the lower positive
equilibrium $E_-$ are saddles. At the upper equilibrium
$E_+$, one has $\Delta_+>0$; hence it is asymptotically
stable when $T_+<0$ and unstable when $T_+>0$.

\paragraph{Test for a Hopf bifurcation.}
The relevant equilibrium for the birth of small limit
cycles is $E_+$. The necessary spectral conditions are
\[
T_+=0,\qquad \Delta_+>0.
\]
The trace condition gives
\begin{equation}
\boxed{
\varepsilon_H=
\frac{2(r+0.1)q_+(1-q_+)}{1+2q_+}.
}
\label{eq:hopf}
\end{equation}
This threshold is positive precisely when
\[
\boxed{r_c<r<\frac23.}
\]
For each fixed $r$ in this interval,
$\Delta_+>0$ at $\varepsilon=\varepsilon_H$, and
\[
\left.
\frac{d\,\operatorname{Re}\lambda}
{d\varepsilon}
\right|_{\varepsilon=\varepsilon_H}
=-\frac{1+2q_+}{2}\neq0.
\]
Thus the spectral and transversality conditions hold.

To conclude that a nondegenerate Hopf bifurcation
actually produces a limit cycle, one must additionally
verify that the first Lyapunov coefficient is nonzero.
Its sign determines the stability of the cycle and
the side of $\varepsilon_H$ on which it exists.
The condition $T_+=0$ alone is not sufficient.

For example, when $r=0.65$,
\[
q_+\approx0.896537216,
\qquad
u_+\approx1.358028411,
\qquad
\varepsilon_H\approx0.049815126.
\]
This is a concrete parameter point at which to perform
the Hopf nondegeneracy test.

\paragraph{A direct nonexistence test.}
For the vector field $(P,Q)$ in \eqref{eq:model},
the Dulac function $B(q,u)=1/q$ on $q>0$ gives
\[
\frac{\partial(BP)}{\partial q}
+\frac{\partial(BQ)}{\partial u}
=
2(r+0.1)(1-q)-\frac{\varepsilon(1+2q)}q.
\]
This expression is nonpositive throughout $q>0$ if
\[
\boxed{
\varepsilon\geq(2-\sqrt3)(r+0.1).
}
\]
The Bendixson--Dulac criterion then excludes all
physical limit cycles. There are also no physical
limit cycles for $r\leq r_c$.

Therefore, the Hopf test identifies where small cycles
can arise, while the Dulac test excludes a region of
parameter space. Neither test by itself determines
the global number of limit cycles.














ppppppppppppppppppppppppp

\section{Preliminary analysis of limit cycles}
\label{sec:preliminaries}

Throughout this section, let
\begin{equation}
A=r+\frac{1}{10},
\qquad
h(q)=\frac{21}{10}-2Aq+Aq^2,
\label{eq:A-h}
\end{equation}
and define the polynomials
\begin{equation}
\begin{aligned}
H(q)
&=(1+2q)h(q)-2-2q\\
&=2Aq^3-3Aq^2+
  \left(\frac{11}{5}-2A\right)q+\frac{1}{10},\\
D(q)
&=\varepsilon(1+2q)+qh'(q)\\
&=2Aq^2+2(\varepsilon-A)q+\varepsilon.
\end{aligned}
\label{eq:H-D}
\end{equation}
The original equations can then be written as
\begin{equation}
\dot q=q[u-h(q)],
\qquad
\dot u=\varepsilon[2+2q-(1+2q)u].
\label{eq:barkley-polynomial}
\end{equation}

\subsection{Invariant region and boundedness}
\label{subsec:invariance}

The half-plane $q\geq0$ is positively invariant because
$q=0$ is an invariant line. Moreover, every solution with
$q(0)>0$ remains in $q>0$ throughout its interval of existence.

The strip
\begin{equation}
\mathcal S=\{(q,u):q\geq0,\ 1\leq u\leq2\}
\label{eq:physical-strip}
\end{equation}
is positively invariant. Indeed,
\[
\left.\dot u\right|_{u=1}=\varepsilon>0,
\qquad
\left.\dot u\right|_{u=2}=-2\varepsilon q\leq0.
\]
More generally, if $q(0)\geq0$, comparison at the horizontal
boundaries gives
\[
\min\{u(0),1\}\leq u(t)\leq\max\{u(0),2\}.
\]
Since $u$ is bounded and the leading term of $\dot q$ is
$-Aq^3$, the component $q$ also remains bounded.
Consequently, every solution starting in $q\geq0$ exists
for all forward time.

Every nonconstant periodic orbit in the physical half-plane
lies in
\begin{equation}
q>0,\qquad 1<u<2.
\label{eq:cycle-strip}
\end{equation}
To see this, first note that a periodic orbit cannot intersect
the invariant line $q=0$, on which
$\dot u=\varepsilon(2-u)$.
At a maximum or minimum of $u$ on a periodic orbit,
\[
u=\frac{2(1+q)}{1+2q}
  =1+\frac{1}{1+2q},
\]
which belongs to $(1,2)$ when $q>0$.


\subsection{Equilibria and their stability}
\label{subsec:equilibria}

The laminar equilibrium
\[
E_0=(0,2)
\]
exists for all $r,\varepsilon>0$. Its Jacobian is
\[
J(E_0)=
\begin{pmatrix}
-\frac{1}{10}&0\\
-2\varepsilon&-\varepsilon
\end{pmatrix},
\]
so $E_0$ is locally asymptotically stable.

A positive equilibrium $E_*=(q_*,u_*)$ satisfies
\begin{equation}
u_*=\frac{2(1+q_*)}{1+2q_*},
\qquad H(q_*)=0.
\label{eq:positive-equilibrium}
\end{equation}
Equivalently, the positive-equilibrium branch is parameterized by
\begin{equation}
r=R(q):=
\frac{1+22q}{10q(2-q)(1+2q)}-\frac{1}{10},
\qquad 0<q<2.
\label{eq:equilibrium-curve}
\end{equation}
Differentiation yields
\begin{equation}
R'(q)=
\frac{C(q)}
{5q^2(q-2)^2(2q+1)^2},
\qquad
C(q)=44q^3-30q^2-3q-1.
\label{eq:R-derivative}
\end{equation}
By Descartes' rule of signs, $C$ has exactly one positive
root. Denote it by $q_{\mathrm{SN}}$. Since
$C(0)<0<C(1)$, this root lies in $(0,1)$.
Thus $R$ decreases on $(0,q_{\mathrm{SN}})$ and increases
on $(q_{\mathrm{SN}},2)$, with
$R(q)\to+\infty$ at both endpoints.

Define
\begin{equation}
r_{\mathrm{SN}}=R(q_{\mathrm{SN}}).
\label{eq:rSN}
\end{equation}
Numerically,
\[
q_{\mathrm{SN}}\approx0.802140182,
\qquad
r_{\mathrm{SN}}\approx0.645189835.
\]
It follows that there are no positive equilibria for
$r<r_{\mathrm{SN}}$, one double positive equilibrium for
$r=r_{\mathrm{SN}}$, and two positive equilibria
\[
E_\pm=
\left(q_\pm,\frac{2(1+q_\pm)}{1+2q_\pm}\right),
\qquad q_-<q_+,
\]
for $r>r_{\mathrm{SN}}$.

At a positive equilibrium, the Jacobian is
\begin{equation}
J_*=
\begin{pmatrix}
2Aq_*(1-q_*)&q_*\\[1mm]
-\dfrac{2\varepsilon}{1+2q_*}&-\varepsilon(1+2q_*)
\end{pmatrix}.
\label{eq:positive-jacobian}
\end{equation}
In particular,
\begin{equation}
\operatorname{tr}J_*=-D(q_*),
\qquad
\det J_*=\varepsilon q_*H'(q_*).
\label{eq:trace-determinant}
\end{equation}
For $r>r_{\mathrm{SN}}$,
\[
H'(q_-)<0,
\qquad H'(q_+)>0.
\]
Hence $E_-$ is a saddle, whereas $E_+$ is asymptotically
stable when $D(q_+)>0$ and unstable when $D(q_+)<0$.

Since $R(1)=2/3$, the upper equilibrium satisfies
\[
q_+<1
\quad\Longleftrightarrow\quad
r_{\mathrm{SN}}<r<\frac23.
\]
For $r\geq2/3$, one has $q_+\geq1$ and therefore
$D(q_+)>0$ for every $\varepsilon>0$.


\subsection{Hopf bifurcation and local limit cycles}
\label{subsec:hopf}

A positive Hopf threshold exists precisely on the upper branch
with $r_{\mathrm{SN}}<r<2/3$. It is given by
\begin{equation}
\varepsilon_H(r)
=
\frac{2Aq_+(1-q_+)}{1+2q_+}.
\label{eq:hopf-threshold}
\end{equation}
At this threshold,
\[
D(q_+)=0,
\qquad
\omega^2:=\varepsilon_Hq_+H'(q_+)>0.
\]
The equilibrium coordinates are independent of $\varepsilon$,
and
\begin{equation}
\left.
\frac{\partial\,\operatorname{Re}\lambda}
{\partial\varepsilon}
\right|_{\varepsilon=\varepsilon_H}
=-\frac{1+2q_+}{2}<0.
\label{eq:hopf-transversality}
\end{equation}
Thus the transversality condition holds.

For the calculation of the Hopf coefficient, the polynomial
substitution
\[
v=u-h(q)
\]
gives
\begin{equation}
\dot q=qv,
\qquad
\dot v=-D(q)v-\varepsilon H(q).
\label{eq:polynomial-auxiliary}
\end{equation}
This auxiliary system remains cubic polynomial and uses
the original time variable.

At the Hopf point, set $s=q_+$ and introduce the linear
coordinates
\[
X=q-s,\qquad Y=-\frac{s}{\omega}v.
\]
The linear part becomes
$\dot X=-\omega Y$, $\dot Y=\omega X$.
The cubic coefficient in the corresponding radial normal
form is
\begin{equation}
\ell=
\frac{D'(s)H''(s)-D''(s)H'(s)}
{16H'(s)},
\label{eq:hopf-coefficient-derivatives}
\end{equation}
where $D$ is evaluated at $\varepsilon=\varepsilon_H$.
Using $H(s)=0$ and \eqref{eq:hopf-threshold}, this simplifies to
\begin{equation}
\ell=
\frac{
(22s+1)
\left(88s^4+176s^3-186s^2+64s+5\right)
}{
80s(2-s)(2s+1)^2C(s)
}.
\label{eq:hopf-coefficient}
\end{equation}
For $q_{\mathrm{SN}}<s<1$, the denominator is positive.
Moreover,
\[
88s^4+176s^3-186s^2+64s+5
=
88s^4+s(176s^2-186s+64)+5>0,
\]
because the quadratic $176s^2-186s+64$ has positive
leading coefficient and negative discriminant.
Consequently,
\begin{equation}
\ell>0.
\label{eq:hopf-positive}
\end{equation}

Therefore, for every fixed
$r\in(r_{\mathrm{SN}},2/3)$, the Hopf bifurcation is
nondegenerate and subcritical. A unique small-amplitude
unstable limit cycle bifurcates from $E_+$ for
$\varepsilon>\varepsilon_H(r)$ sufficiently close to
$\varepsilon_H(r)$. The equilibrium inside this cycle
is asymptotically stable.

This uniqueness is local: it concerns cycles sufficiently
close to the Hopf equilibrium and does not exclude additional
cycles elsewhere in the phase plane.


\subsection{Restrictions on the location of periodic orbits}
\label{subsec:cycle-location}

At an extremum of $q$ on a nonconstant periodic orbit,
$\dot q=0$ and hence $u=h(q)$. Differentiating the first
equation of \eqref{eq:barkley-polynomial} gives, at such a point,
\begin{equation}
\ddot q=-\varepsilon qH(q).
\label{eq:q-extremum}
\end{equation}
The extremum cannot satisfy $H(q)=0$, since it would then
be an equilibrium point.

If $r\leq r_{\mathrm{SN}}$, one has $H(q)\geq0$ for all
$q>0$. Equation~\eqref{eq:q-extremum} then excludes the
minimum required by any nonconstant periodic orbit.
Thus
\begin{equation}
0<r\leq r_{\mathrm{SN}}
\quad\Longrightarrow\quad
\text{no limit cycles in }q\geq0.
\label{eq:no-cycles-small-r}
\end{equation}

Suppose now that $r>r_{\mathrm{SN}}$. Since $H$ is negative
precisely between its positive roots, the minimum and maximum
of $q$ along any periodic orbit satisfy
\begin{equation}
q_-<q_{\min}<q_+<q_{\max}<2.
\label{eq:cycle-location}
\end{equation}
For the final inequality, recall that $u<2$ along the orbit,
whereas
\[
h(q)=\frac{21}{10}+Aq(q-2)\geq\frac{21}{10}
\qquad\text{for }q\geq2.
\]
Hence a $q$-maximum cannot occur at $q\geq2$.

In particular, every physical limit cycle excludes the saddle
$E_-$ and the boundary equilibrium $E_0$. By the index
theorem for planar periodic orbits, it surrounds $E_+$.


\subsection{A Bendixson--Dulac exclusion region}
\label{subsec:dulac}

Let $P$ and $Q$ denote the right-hand sides of
\eqref{eq:barkley}. On the simply connected half-plane
$q>0$, choose the Dulac function
\[
B(q,u)=\frac1q.
\]
A direct calculation gives
\begin{equation}
\frac{\partial(BP)}{\partial q}
+
\frac{\partial(BQ)}{\partial u}
=
2A(1-q)-\frac{\varepsilon(1+2q)}q
=
-\frac{D(q)}q.
\label{eq:dulac-divergence}
\end{equation}
The quadratic polynomial $D$ is nonnegative for all $q>0$
if
\begin{equation}
\varepsilon\geq(2-\sqrt3)A.
\label{eq:dulac-condition}
\end{equation}
Indeed, for $\varepsilon\geq A$ all its coefficients are
nonnegative. For $0<\varepsilon<A$, its minimum is
\[
\min_{q>0}D(q)
=
\varepsilon-\frac{(A-\varepsilon)^2}{2A},
\]
which is nonnegative precisely when
$\varepsilon\geq(2-\sqrt3)A$.

The Bendixson--Dulac criterion therefore yields
\begin{equation}
\boxed{
\varepsilon\geq
(2-\sqrt3)\left(r+\frac1{10}\right)
\quad\Longrightarrow\quad
\text{no limit cycles in }q\geq0.
}
\label{eq:no-cycles-large-epsilon}
\end{equation}
At equality, the Dulac divergence vanishes only on one
vertical line, which does not affect the conclusion.


\subsection{Consequences for the global cycle problem}
\label{subsec:global-cycle-problem}

The preceding results restrict the search for physical
limit cycles to
\[
r>r_{\mathrm{SN}},
\qquad
0<\varepsilon<
(2-\sqrt3)\left(r+\frac1{10}\right).
\]
Every such cycle surrounds the upper positive equilibrium
and remains in the region specified by
\eqref{eq:cycle-strip} and \eqref{eq:cycle-location}.
Furthermore, a small unstable cycle is known to exist on
the subcritical side of the Hopf curve
\eqref{eq:hopf-threshold}.

These conclusions establish local existence and explicit
global exclusion regions. They do not yet determine the
total number of limit cycles or the global continuation
of the Hopf branch. In particular, uniqueness away from
the Hopf point and possible termination in a saddle
homoclinic loop require additional analysis.







%\section*{Acknowledgments}
%[Funding, collaborations]

\bibliographystyle{plain}
\bibliography{bibliography}
%\printbibliography




\end{document}
